% =====================================================================
% Related Literature Review, v5
% DRAFT (Claude, 2026-11-09, with Anders). Rewrite of v4 (the draft written
% without access to the papers). Every claim below is checked against the
% paper notes in literature/papers/<bibkey>/note.md (not against paper.txt
% directly; corrected 2026-11-10). The corrections made to v4 are listed in
% drafts/literature-review-v5/CHANGES.md. Text for the students to rework.
% Structure agreed with Anders 2026-11-09: scope, valuation and capacities,
% uncertainty and price dynamics, market participation, solution methods,
% research gap.
% Citations: natbib commands (\citet, \citep) through biblatex natbib=true.
% Keys from ../pumped-hydro-thesis/references.bib.
% =====================================================================

\section{Related Literature Review}
\label{sec:literature}

Studies of pumped-storage hydropower differ in what they value, which
uncertainties they represent, which markets they include and how they
solve the resulting planning problem. This chapter reviews these four
themes in turn. Section~\ref{sec:lit_scope} describes how the review is
limited and how the literature was found.
Section~\ref{sec:lit_valuation} reviews how the value of pumped storage
and of its capacities has been assessed, and
Section~\ref{sec:lit_uncertainty} how uncertainty and price dynamics have
been represented. Section~\ref{sec:lit_markets} reviews participation in
energy and reserve markets, and Section~\ref{sec:lit_methods} the
solution methods used and the physical detail they allow. Finally,
Section~\ref{sec:lit_gap} identifies the research gap and positions this
study. Table~\ref{tab:lit_comparison} compares the reviewed studies.

% ---------------------------------------------------------------------
\subsection{Scope and Search}
\label{sec:lit_scope}

The review covers the operational scheduling and economic valuation of
pumped-storage hydropower, with emphasis on studies that link short-term
operation to the long-term management of stored water. Studies of
conventional reservoir hydropower are included where they show how
uncertainty, reserve markets or plant detail affect water values, since a
reservoir plant faces the same storage decision without the option to
pump. System studies and studies of hybrid island systems are included
where they assess pumped-storage capacity, because they show how the
value of a capacity depends on the objective of the study. Studies of
price dynamics and storage value are included because the research
question concerns price volatility and persistence.

\todonote{TODO (authors): describe the search you actually carried out,
in three to five sentences: databases (e.g.\ Google Scholar, Scopus,
Oria), search strings, inclusion criteria (peer-reviewed, title and
abstract within scope) and whether further studies were found by
following reference lists (backward snowballing, as Bore and Brensholm
describe in their Section~4.1). The search log
(\texttt{literature/search-log.md}) is still empty, so nothing may be
stated here before it has been logged.}

% ---------------------------------------------------------------------
\subsection{Valuation of Pumped Storage and Its Capacities}
\label{sec:lit_valuation}

The value of pumped storage depends on whose value is measured. System
studies measure the change in total surplus, whereas owner studies
measure the operating profit of one plant. The two can point in opposite
directions. \citet{henden2016economic} add between 0 and
\SI{14250}{\mega\watt} of new pumped storage in southern Norway to a
model of the Northern European power system in 2030. Only about
\SI{950}{\mega\watt} maximises socio-economic surplus, and owners need
total exchange capacity of more than about \SI{11000}{\mega\watt}
before the investment pays off. With high transmission capacity, a large pumped-storage
fleet runs often but earns almost nothing, because its own operation
removes the price differences it trades on. \citet{dallinger2019socioeconomic}
find the same effect in Central Europe for 2030: the average price spread
captured by pumped storage falls as Austrian storage capacity grows.
These studies show that the value of storage is not a fixed property of
the plant but falls as more storage is built. They value energy trading
only, which the authors of the first study expect to understate the
value of pumped storage. For a single owner, the result means that the
price-taker assumption holds only if the plant is small relative to the
market it trades in.

\citet{helseth2017} assess an expansion of the \SI{200}{\mega\watt} Duge
pumped-storage plant in southern Norway by up to \SI{2000}{\mega\watt}
in a detailed model of the Nordic system. Their detailed model gives a
smoother, but lower, marginal value of the additional capacity than an
aggregated model, because it represents the constraints of the
surrounding watercourse. The paper reports the result as a change in
socio-economic surplus and does not state whether pump and turbine
capacity are expanded separately. It therefore shows that the
representation of the plant changes the value of capacity, but not which
capacity is worth expanding.

Owner studies that add a pump to an existing cascade report small
gains. \citet{ak2019quantifying} convert a three-reservoir Turkish
cascade into a pumped-storage system and find that net revenue rises by
2.9 to 10.4\%, 6.0\% on average, across five historical price years.
Turbine revenue rises by 23 to 32\%, but most of this gain is offset by
the cost of pumping. \citet{alic2023optimal} retrofit a \SI{38}{\mega\watt}
pump to the upper plant of a Norwegian cascade and find a median revenue
gain of only 2.35\%. \citet{schuller2024practical} value a new
\SI{90}{\mega\watt} plant with pumping in the Sira-Kvina system and find
that net income rises by 2.9\% in their medium-term model but by 1.9\%
when the operation is simulated with hourly detail, because short-term
restrictions, such as a minimum head for pumping, reduce the use of the
new plant. All three studies use the energy market only, and each
evaluates a single pump size. In \citet{alic2023optimal}, the pump also
has the same power as the turbine, so the effect of pump size cannot be
separated from that of adding a pump at all.

Studies that include reserve markets find larger and more varied gains
from a change in configuration. \citet{helseth2017iet} compare two
reinvestment alternatives of \SI{24.8}{\mega\watt} for a Norwegian
reservoir plant. Selling primary reserves raises expected operating
profit by 7.3\% for a configuration with a large Francis and a small Pelton
unit and by 10.8\% for one with two Pelton units, since Pelton units can run at low
output and leave more capacity for reserves. \citet{chazarra2017}
compare a conventional closed-loop plant with two plants that can pump
and generate at the same time. The configuration with ternary units
earns 141 to 166\% more than the conventional plant in two
representative weeks, whereas the configuration with reversible units
earns only 2.8 to 3.5\% more. In both studies, the configuration that
gives the most flexible operating range wins because reserve capacity is
paid. Neither varies the size of the reservoir. \citet{chazarra2017}
note that the flexible plant uses only a small part of its storage and
suggest, without testing it, that this could affect reservoir sizing.

\citet{bore2026assessing} come closest to valuing capacity for an owner
under uncertainty. They add 30, 60 and
\SI{90}{\mega\watt} of generating capacity to a conventional plant in
southern Norway. Selling reserve capacity raises revenue by 25.8\% compared with energy
trading alone for the plant without an upgrade, and by 57.5\% with the
\SI{90}{\mega\watt} upgrade \citep{bore2026}. Energy
revenue hardly changes with the upgrade, because the plant has no more
water to sell. The extra capacity is therefore paid mainly through
reserve markets, and the marginal benefit of each step falls
\citep{bore2026assessing}. Their upgrades change only generating power,
and the plant has no pump.

The value of reservoir capacity has been examined mainly outside
electricity markets. \citet{zisos2023mixing} size the upper reservoir of
a seawater pumped-storage plant on a Greek island, with pump and turbine
both fixed at \SI{6}{\mega\watt}. The optimal storage is
\SI{315195}{\cubic\metre} for the deterministic case and about
\SI{330000}{\cubic\metre} on average over 100 weather and demand
scenarios, less than a third of the
\SI{1.1}{\mega\cubic\metre} of an earlier design. The plant follows a
fixed rule that pumps all surplus and covers all deficits, and its value
comes from reduced unmet demand rather than from prices. The study shows
that the reservoir size matters, but with fixed pump and turbine power
and without prices it cannot show how the three capacities interact.

The reviewed studies thus vary one capacity at a time: a pump is added,
a generator is upgraded or replaced, or a reservoir is sized. None
varies pumping power, generating power and reservoir capacity jointly
for a plant that trades in markets. This matters because one capacity
may limit another. The finding of \citet{bore2026assessing} that extra
generating power earns little energy revenue when water is scarce is one
example; a larger pump adds little when the reservoir is already full.

% ---------------------------------------------------------------------
\subsection{Uncertainty and Price Dynamics}
\label{sec:lit_uncertainty}

Prices and inflow are unknown when the owner decides whether to store
water, and the value of storage depends on how they develop. The
reviewed studies differ in which inputs they treat as uncertain, how
they describe them and how much the decision maker is assumed to know.

Several studies describe the future through historical or simulated
years and optimise each year with full knowledge of it.
\citet{ak2019quantifying} solve the cascade problem for each of five
historical price years, so each result assumes that the whole year is
known in advance. \citet{chazarra2017} schedule each day with known
prices and reserve activation, and \citet{naval2025optimization} treat
the 2023 Spanish day-ahead prices as known for a whole year while
treating only wind as uncertain. The reported revenues are therefore
upper bounds on what an operator can earn, and the authors of the
second study state this themselves. \citet{helseth2017} let each week's
decision be informed by 19 weather and price scenarios drawn from
historical years. All of these scenarios are revealed at once after the
first week, so the plan for later weeks is made with knowledge of each
scenario.

Other studies let information arrive week by week. \citet{alic2023optimal}
cluster weekly price and inflow observations into five joint states per
week and estimate the probabilities of moving between them.
\citet{helseth2017iet} and \citet{helseth2016} build weekly states for
energy and reserve prices, and \citet{helseth2016} model inflow
separately. \citet{bore2026assessing} use three joint price and inflow
states per week. Such a description keeps the dependence between
consecutive weeks, but each week's outcome is represented by a handful
of typical weeks. In \citet{alic2023optimal}, weekly average prices are
spread over the week with a fixed profile, so the variation of prices
within the week, which a pump exploits, is the same in every state. This
may explain part of their small gain from pumping, although the authors
do not say so.

How the dependence between prices and inflow is described appears to
matter less for revenue than for operation. \citet{kleiven2023} find
that ignoring the negative correlation between Norwegian prices and
inflow lowers water values and leads to a higher reservoir in autumn,
but costs only 0.17 to 0.30\% of expected revenue. \citet{perezdiaz2020medium}
likewise find that the correlation between energy and reserve prices has
little effect on water values for a Spanish reservoir plant. These
results concern conventional plants and say little about whether the
same holds when prices also drive pumping.

The information available within the week is a separate question.
\citet{abgottspon2014} compare four ways of describing the week for a
Swiss pumped-storage plant. Assuming that the whole week is known
overvalues water. Representing the week as uncertain raises expected
profit and its conditional value at risk by about 10\% when weekly
secondary reserves are included, but gives no advantage without them.
\citet{alic2023optimal} show a related effect of time resolution: when
operation is simulated day by day instead of week by week, revenue falls
by 24\% in the first week in which an environmental restriction on the
reservoir applies. Both results suggest that a known week is a
reasonable simplification for energy trading, but less so when reserves
or state-dependent restrictions are included.

Few of the reviewed studies examine how the behaviour of prices over
time affects the value of storage, although this is what a pump and a
reservoir earn from. Electricity prices in the Nordic market are both
volatile and persistent. \citet{haldrup2006regime} find that hourly
Nord Pool prices in 2000--2003 revert to their mean only slowly, with
the strongest persistence in the Norwegian price areas, which they link
to the long memory of reservoirs and inflow. Regimes with and without
transmission congestion last on average from about 5 to 77 hours.
\citet{mo2001} note that a description in which next week's price
depends only on this week's state forgets the past quickly, so price
uncertainty about prices far into the future needs an additional model.
Section~\ref{sec:ph25-volatility-persistence} defines volatility and
persistence as used in this study.

The link between price variation and storage value has mainly been
examined with known prices. \citet{ak2019quantifying} find the largest
gain from pumping in the price year with the largest variation within
the day, but the relation is not one-to-one, because changes from month
to month and head also matter. \citet{sioshansi2009estimating} find
that, with known prices in the PJM market, the first four hours of
storage earn more than half of the arbitrage value, eight hours about
85\% and twenty hours about 95\%. \citet{bradbury2014economic} show,
for seven US markets, that greater price volatility raises the return on
storage, but that the most profitable ratio of storage to power is set
mainly by efficiency and by how the cost is split between power and
storage, not by volatility. \citet{lohndorf2023} report that
coordinating day-ahead and intraday trading is worth 10.0 to 29.1\% for
a pumped-storage plant but only 1.0 to 1.6\% for a reservoir plant
without a pump, which suggests that short-term price variation matters
more for a plant that can pump. None of these studies varies volatility
and persistence separately, and none relates them to the value of pump,
turbine and reservoir capacity under uncertainty. Since a pump and a
turbine earn from short price differences, whereas a reservoir earns
from differences that persist over weeks or seasons, the two properties
may favour different capacities.

% ---------------------------------------------------------------------
\subsection{Market Participation}
\label{sec:lit_markets}

The markets included decide which revenues a study can report. Several
of the reviewed owner studies include only day-ahead energy
\citep{ak2019quantifying,alic2023optimal,schuller2024practical}, as do
the system studies of \citet{henden2016economic} and \citet{helseth2017}.
\citet{abgottspon2013stochastic} add forward contracts for a Swiss plant
with \SI{240}{\mega\watt} of generating power, \SI{46}{\mega\watt} of
pumping power and about \SI{200}{\giga\watt\hour} of storage, which they
model as a producer that affects prices. Forward trading raises
simulated annual profit from about 20 to 24 million euro. Forward
contracts spread the sales of the same energy over time, however, and
do not pay for flexibility as reserves do.

For conventional reservoir plants, reserve sales consistently raise
profit and change how the reservoir is used. \citet{helseth2016} sell
symmetric reserve capacity, traded the week before energy, from a
Norwegian watercourse of seven plants. A large part of the reserve
profit comes without changing the generation schedule, because the
plants are most efficient below full output and therefore have spare
capacity both up and down. With reserves, the large reservoirs are kept
higher until the spring flood and then used more to keep generators
running. \citet{helseth2017iet} find the same pattern for a single
reservoir, together with water values that fall faster as the reservoir
fills, because capacity held for reserves limits how fast water can be
discharged. \citet{perezdiaz2020medium} find that including the
secondary reserve market when computing water values raises the profit
of a Spanish reservoir plant by about 4.5\%.
\citet{bore2026assessing} add intraday energy and four Nordic reserve
products and find that reserve revenues provide a hedge, since they make
up a larger share of revenue in periods with high inflow and low energy
prices.

A pump extends what a plant can offer in reserve markets.
\citet{dallinger2019socioeconomic} show that adding pump-turbines to an
Austrian reservoir plant extends its operating range from minimum load
up to full output to between about $-70$\% and 100\% of turbine
capacity. The plant can then offer downward reserve without generating,
and its ability to deliver reserves no longer depends on natural inflow.
\citet{chazarra2017} find that a closed-loop plant earns most of its
income from secondary reserve capacity, and that the day-ahead market
becomes a means of enabling reserve sales. A conventional plant cannot
offer reserves while pumping, which is why the configuration that pumps
and generates at the same time earns much more. \citet{abgottspon2014}
include weekly secondary reserves for a Swiss pumped-storage plant and
find that they raise expected profit for all four models compared.

The treatment of reserve activation differs widely. \citet{chazarra2017}
include the energy delivered when reserves are activated, based on
historical activation, and show that storage limits restrict reserve
offers: the conventional plant does not offer reserves in an hour when
the expected activation of about 85\% could empty the upper reservoir.
\citet{dallinger2019socioeconomic} simulate activations every quarter
hour, but offer capacity at a price of zero, so their reserve revenue
comes from activated energy only. The other reviewed studies pay for
capacity and leave activation out
\citep{helseth2016,helseth2017iet,abgottspon2014,bore2026assessing}.
Leaving activation out omits its energy settlement and its effect on
stored water, which matters most for a plant with a small reservoir.
Most evidence on reserve markets and water values concerns conventional
plants. For pumped storage, only \citet{abgottspon2014} combine reserves
with water values over a year, for a single plant whose capacities are
not reported.

% ---------------------------------------------------------------------
\subsection{Solution Methods}
\label{sec:lit_methods}

The studies above combine a long horizon, uncertain inputs, several
markets and different levels of plant detail. No method handles all of
these equally well, and the physical detail a study can keep depends on
the method it uses. This section reviews the methods applied, whose
principles are described in Sections~\ref{subsec:two_stage_stochastic_optimisation},
\ref{subsec:stochastic_dynamic_programming} and
\ref{subsec:stochastic_dual_dynamic_programming}.

\subsubsection*{Deterministic optimisation and simulation}

Deterministic models solve one problem with given inputs and can
therefore include nonlinear and discrete plant features.
\citet{chazarra2017} use a mixed-integer linear model to represent unit
modes and simultaneous pumping and generation, solved one day at a time
with the same reservoir level at the start and end of each day.
\citet{ak2019quantifying} use a nonlinear model with head-dependent
generation and monthly decisions, in which hourly prices enter only
through representative price curves. \citet{naval2025optimization} solve
a mixed-integer model for a hybrid wind, solar and pumped-storage system
over a year, with chance constraints on how much of the demand renewables
cover, and \citet{zisos2023mixing} combine simulation with a heuristic
search over reservoir sizes. These models give results with full
knowledge of the inputs, so they do not provide a value of stored water
that an operator could use without knowing the future. Where the
horizon is longer than a day, the use of the reservoir is fixed by end
conditions or monthly rules rather than by a water value.

\subsubsection*{Two-stage stochastic programming in a rolling horizon}

\citet{helseth2017} solve, every week, a two-stage stochastic linear
programme in which the first stage is the current week and the second
stage contains one scenario for each of 19 historical years, and then
move one week forward. This keeps all 1265 hydropower modules of the
Nordic system without aggregation and uses historical weather directly.
The price is a long run time, 50 hours for a transmission case and
156 hours for the pumped-storage case, and the plan for later weeks is
made with knowledge of each scenario. The value of water at the end of
the scenarios is taken from another model.

\subsubsection*{Stochastic dynamic programming}

Stochastic dynamic programming computes water values on a grid of
reservoir levels for each stage and state, and places no convexity
requirement on the stage problems. \citet{alic2023optimal} choose it
because their environmental restrictions depend on the reservoir level
and are not convex, even though they note that stochastic dual dynamic
programming is more efficient for medium- and long-term scheduling.
\citet{helseth2017iet} solve weekly mixed-integer problems with
head-dependent generation and unit commitment. A linear version of their
model overestimates profit with reserve sales by 2.3 to 5.2 percentage
points, mainly because it lets the plant sell reserves below its minimum
output. Water values from the linear model, used in a simulation with
the mixed-integer model, lose 0.6 to 1.0 percentage points. The cost of
this detail is the grid: \citet{alic2023optimal} need hours to solve two
reservoirs with 25 grid points each, and the reviewed applications
concern one plant or a small cascade.

\subsubsection*{Stochastic dual dynamic programming}

Stochastic dual dynamic programming \citep{pereira1991} avoids the grid
by approximating the expected future profit with linear cuts that are
refined iteratively, and has become the standard method for medium- and
long-term hydropower scheduling \citep{fullnerRebennack2025}. Its
convergence results require convex, typically linear, stage problems
\citep{shapiro2011}. \citet{helseth2013} apply it to the Icelandic
hydrothermal system with a \SI{100}{\mega\watt} pumped-storage plant,
and describe the exogenous price by weekly states between which the
cuts are computed separately, which keeps the problem convex. Their
linear approximation of pump start-up costs reduces switching from
generation to pumping from 6.54 to 0.12\% of the time periods, but the
pump still runs below its minimum output, so the approximation is
imperfect. \citet{helseth2016} use the same combination of price states
and cuts to include uncertain reserve prices, and note that the linear
plant description may cause significant errors when reserves are sold.

Several studies relax the linearity. \citet{hjelmeland2019energies}
include head-dependent generation and use strengthened cuts. These cuts
give a better operating policy than ordinary cuts at about five times
the run time, but with uncertain prices the gap between the bounds is
still 36\%. \citet{helseth2024approximating} represent the detailed
Norwegian hydropower system in an aggregated model through feasibility
cuts computed from the detailed system, building on
\citet{helseth2023hydropowera}. On 30 historical years this lowers total
cost by 7.2\% on average compared with plain aggregation, at 2.7 times
the run time. These extensions bring the plant description closer to
that of stochastic dynamic programming, but at a computational cost.

The method has also been used for single pumped-storage plants.
\citet{abgottspon2016multihorizon} schedule a Swiss system of ten
reservoirs, with \SI{1125}{\mega\watt} of turbines and
\SI{424}{\mega\watt} of pumps, by using only the four large reservoirs
as states between days and keeping the six small reservoirs inside
hourly problems within each day. Their model reaches a gap of 0.5\%, but the
reported case is deterministic. \citet{schuller2024practical} use water
values from a medium-term model of this type to value a pumping
investment. \citet{bore2026assessing} use the method with three price
and inflow states per week for a conventional plant. In their model,
each iteration samples a single path through the weeks, the algorithm
stops after a fixed 50 iterations, and no bounds are reported. They also
find that the water value does not fall as the reservoir approaches
full, which they attribute to the few reservoir levels the algorithm
visits.

The methods thus trade plant detail against the size of the problem and
the description of uncertainty. Deterministic models keep the most
detail but give no water values under uncertainty, and two-stage models
add scenarios but reveal them all at once. Stochastic dynamic
programming keeps detail and sequential decisions for small systems,
whereas stochastic dual dynamic programming handles more states and
finer resolution if the stage problems are linear. A study that
compares many configurations under several price cases needs a method
whose run time allows many solutions, and the results of
\citet{helseth2017iet} indicate what a linear plant description costs.

% ---------------------------------------------------------------------
\subsection{Research Gap and Positioning of This Study}
\label{sec:lit_gap}

Table~\ref{tab:lit_comparison} summarises the reviewed studies. The
literature shows that reserve markets change both the profit and the
water values of a hydropower plant, that pumping extends what a plant
can offer in these markets, and that the value of storage falls as more
storage is built. It also shows that the gains from adding a pump to an
energy-only plant are small, from about 2 to 10\%, while flexible
configurations in reserve markets can earn much more.

The closest studies each cover part of the question of which capacities
create value. \citet{bore2026assessing} value generating-capacity
upgrades under uncertainty in several markets, but for a plant without a
pump. \citet{alic2023optimal} and \citet{schuller2024practical} value a
single pump in the energy market, \citet{helseth2017iet} and
\citet{chazarra2017} compare unit configurations of given size, and
\citet{zisos2023mixing} size a reservoir with fixed pump and turbine and
no market. \citet{helseth2017} expand a pumped-storage plant from a
system perspective without separating its capacities. No reviewed study
varies pumping power, generating power and reservoir capacity jointly
for a price-taking owner in energy and reserve markets under uncertainty.
In addition, the effect of price behaviour on storage value has been
examined only with known prices, and never with volatility and
persistence varied separately.

This study therefore asks how pumping power, generating power and
reservoir capacity, individually and in combination, affect the
long-term operating value of a price-taking owner that trades in
day-ahead and reserve-capacity markets, and how the ranking of these
capacities changes with price volatility and persistence. Following the
reviewed applications to reservoir hydropower, the problem is solved
with stochastic dual dynamic programming, with weekly stages and
uncertain inflow and day-ahead prices. The last row of
Table~\ref{tab:lit_comparison} places the study among the reviewed
literature. The contribution is a comparison of capacity combinations
under uncertainty, not a claim that stochastic scheduling, reserve
participation or storage sizing is new on its own.

% ---------------------------------------------------------------------
% Table 3. Check the row "This study" against ch. 6 when it changes
% (reserve products depend on Q-002).
\begin{landscape}
\begingroup
\footnotesize
\setlength{\tabcolsep}{4pt}
\renewcommand{\arraystretch}{1.15}
\begin{longtable}{@{}>{\raggedright\arraybackslash}p{0.12\linewidth}
                   >{\raggedright\arraybackslash}p{0.15\linewidth}
                   >{\raggedright\arraybackslash}p{0.19\linewidth}
                   >{\raggedright\arraybackslash}p{0.17\linewidth}
                   >{\raggedright\arraybackslash}p{0.18\linewidth}
                   >{\raggedright\arraybackslash}p{0.13\linewidth}@{}}
\caption{Reviewed studies compared by perspective, uncertainty, markets,
capacity question and solution method, and the position of this study.}
\label{tab:lit_comparison}\\
\toprule
\textbf{Study} & \textbf{Perspective and asset} & \textbf{Uncertainty and information} & \textbf{Markets} & \textbf{Capacity question} & \textbf{Solution method}\\
\midrule
\endfirsthead
\multicolumn{6}{l}{\tablename~\thetable\ (continued)}\\
\toprule
\textbf{Study} & \textbf{Perspective and asset} & \textbf{Uncertainty and information} & \textbf{Markets} & \textbf{Capacity question} & \textbf{Solution method}\\
\midrule
\endhead
\midrule
\multicolumn{6}{r}{\footnotesize Continued on the next page}\\
\endfoot
\bottomrule
\addlinespace
\multicolumn{6}{@{}p{\linewidth}@{}}{\footnotesize SDP: stochastic
dynamic programming; SDDP: stochastic dual dynamic programming; LP:
linear programming; MILP: mixed-integer linear programming; NLP:
nonlinear programming; PHS: pumped-storage hydropower; DA: day-ahead; ID: intraday; FCR-N/FCR-D:
frequency containment reserve for normal operation/disturbances; aFRR:
automatic frequency restoration reserve. Studies cited only for price
dynamics or as general methodological references are not included.}\\
\endlastfoot
\citet{helseth2013} & System; Icelandic hydrothermal system with a \SI{100}{\mega\watt} PHS & Inflow, wind and exogenous price; weekly stages, week known & No market products; exogenous price & Operation; effect of grid and pump start-up costs & SDDP with price states\\
\citet{abgottspon2013stochastic} & Price-making owner; Swiss PHS, 240/\SI{46}{\mega\watt}, \SI{200}{\giga\watt\hour} & Inflow; one year & DA and forward contracts & Operation of the given plant & SDP\\
\citet{abgottspon2014} & Owner; Swiss PHS, capacities not reported & Inflow and prices; four descriptions of the week compared & DA and weekly secondary reserve capacity & Operation of the given plant & SDP with hourly within-week problems\\
\citet{helseth2016} & Price-taking owner; Norwegian watercourse, seven plants & Inflow, energy and reserve prices; weekly stages & DA and reserve capacity traded the week before & Operation; water values with reserves & SDDP with price states\\
\citet{helseth2017iet} & Price-taking owner; Norwegian reservoir plant, no pump & Inflow, energy and reserve prices; weekly states & DA and primary reserve capacity & Two reinvestment alternatives of \SI{24.8}{\mega\watt} & SDP with weekly MILP\\
\citet{chazarra2017} & Price-taking owner; closed-loop PHS & None; known prices and activation, daily & DA, secondary reserve capacity and activation energy & Three unit configurations, simultaneous pumping and generation & MILP\\
\citet{helseth2017} & System; Nordic power system & Weather and exogenous prices; all scenarios revealed after the first week & No market products & Duge PHS expanded by up to \SI{2000}{\mega\watt} & Two-stage stochastic LP in a rolling horizon\\
\citet{ak2019quantifying} & Price-taking owner; Turkish cascade & None; each historical price year known & DA energy & Existing cascade versus pumped-storage operation & NLP with monthly decisions\\
\citet{dallinger2019socioeconomic} & System and owner; Central Europe and an Austrian plant & None; one historical year known & DA and aFRR (activation energy only) & Adding pump-turbines; lower minimum load & Dispatch model and MILP\\
\citet{alic2023optimal} & Price-taking owner; Norwegian cascade & Inflow and price; weekly states & Energy & Adding a \SI{38}{\mega\watt} pump & SDP\\
\citet{zisos2023mixing} & Island system; seawater PHS & Wind, demand and power curve; 100 scenarios & None & Reservoir and solar sizing; pump and turbine fixed at \SI{6}{\mega\watt} & Simulation with heuristic search\\
\citet{schuller2024practical} & Price-taking owner; Sira-Kvina system & Inflow and price & Energy & A new \SI{90}{\mega\watt} plant with pumping & SDDP-based water values, hourly simulation\\
\citet{naval2025optimization} & Owner of a hybrid system; closed-loop PHS & Wind; prices known for the year & DA energy & Operation of the given system & MILP with chance constraints\\
\citet{bore2026assessing} & Price-taking owner; Norwegian reservoir plant, no pump & Inflow, DA price; capacity prices at three levels; week known & DA, ID, FCR-N, FCR-D and aFRR capacity; activation excluded & Generating-power upgrades of 30, 60 and \SI{90}{\mega\watt} & SDDP with price and inflow states\\
\midrule
This study & Price-taking owner; pumped storage with natural inflow & Inflow and DA price; weekly stages, week known; varied price volatility and persistence & DA and reserve capacity; activation excluded & Combinations of pumping power, generating power and reservoir capacity & SDDP with price and inflow states\\
\end{longtable}
\endgroup
\end{landscape}
